\documentclass[10pt,a4paper]{amsart}
\usepackage[margin=19mm]{geometry}
\usepackage{amsmath,amssymb,mathtools}
\usepackage[scr=rsfs,cal=euler]{mathalfa}
\usepackage{xfrac}
\usepackage[unicode,hidelinks,bookmarks=false]{hyperref}
\newcommand{\ie}{i.e.\ }
\newcommand{\cf}{cf.\ }
\newcommand{\resp}{respectively\ }
\newcommand{\Subsection}{\subsection}
\providecommand{\nobreakdash}{\nobreak\hbox{-}}
\setlength{\emergencystretch}{2em}
\multlinegap0pt
\usepackage{bm}
\usepackage{bbm}
\usepackage{dsfont}
\usepackage{xspace}
\usepackage{cancel}
\usepackage{mathtools}
\usepackage{amsthm}
\makeatletter
\@ifundefined{theorem}{\newtheorem{theorem}{Theorem}}{}
\@ifundefined{definition}{\newtheorem{definition}[theorem]{Definition}}{}
\makeatother
\makeatletter
\expandafter\ifx\csname rhp\endcsname\relax
\newenvironment{rhp}[1][]{\refstepcounter{rhp}\par\medskip
   \noindent \textbf{Riemann-Hilbert problem~\therhp. #1} \rmfamily}{\medskip}
\fi
\makeatother
\title[Recurrence relations and the Christoffel-Darboux formula for]{Recurrence relations and the Christoffel-Darboux formula for elliptic orthogonal polynomials}
\author{Harini Desiraju, Sampad Lahiry}
\date{Source: arXiv:2506.09582v1 (CC BY 4.0)}
\begin{document}
\maketitle

\noindent\emph{Source and licence.} Recurrence relations and the Christoffel-Darboux formula for elliptic orthogonal polynomials. Version arXiv:2506.09582v1 (CC BY 4.0), distributed under the Creative Commons Attribution 4.0 International licence, as recorded in the arXiv licence metadata for this exact version and shown on the abstract page https://arxiv.org/abs/2506.09582v1. Full rights notice: licenses/arxiv-2506.09582-CC-BY-4.0.txt.\par\smallskip

\noindent\textit{Editorial scope.} Counted unit: the Theorem whose source label is \texttt{5t} in the article (source0.tex lines 452--471, source label \texttt{5t}), delivered together with the article's own complete original proof of it (source0.tex lines 472--524, one \texttt{proof} environment of 53 lines). No mathematics is written, repaired or re-derived: statement and proof are the article's own text line for line. Required context retained uncounted: def:orthonormalEOPs (lines 263--265). Every definition, display and label of those lines is retained unchanged. Omitted from this package and not redistributed: 1--262, 266--451, 525--1215. Nothing inside a retained excerpt is omitted or altered: every retained body is delivered exactly as the article's lines are written, so no omission receipt of any kind accompanies this package. Citations: the retained excerpts carry no citation command at all, so no bibliography entry of the article is transcribed and no reference is redistributed. Reference adaptation: none was needed and none was made; every reference inside the delivered unit resolves inside it, so no unresolved reference or citation is produced. Printed slips kept literal: no printed word, symbol, hypothesis or number of the counted statement or of its proof was altered. Layout and class: the article's own class is \texttt{article} with packages including color, soul, amsmath, amssymb, amsthm, latexsym, babel, graphicx, enumitem, comment, authblk, csquotes, circuitikz, tikz; the wrapper is the lane's pinned \texttt{amsart} editorial wrapper at 10pt on A4, which restates the theorem-like environments the retained text needs and adds the article's own macro definitions that the retained text needs (none), with extra wrapper packages (bm, bbm, dsfont, xspace, cancel, mathtools). The delivered numbering is the unit's own. Assets: no retained span contains a figure environment, a graphics-inclusion command, a PDF-page inclusion, a table or a TikZ picture; no image file is copied into this package, the package declares no asset dependency and no image file is redistributed with it. Licence: version arXiv:2506.09582v1 of this article is distributed under the Creative Commons Attribution 4.0 International licence, as recorded in the arXiv licence metadata for this exact version and shown on the abstract page https://arxiv.org/abs/2506.09582v1. A full rights notice is delivered at licenses/arxiv-2506.09582-CC-BY-4.0.txt. Characters: every retained excerpt is pure ASCII, so the delivered engine, XeLaTeX with its default Latin Modern text font, reports no missing character.
\begin{definition}\label{def:orthonormalEOPs}
 The family of EOPs \( \{\pi_n(z)\}_{n \geq 0,\, n \neq 1} \) are orthonormal with respect to the weight function \( w(z) \) on the contour \( \gamma \).
\end{definition}

\begin{theorem}
Let \( \{a_n>0\}, \{b_n\}, \{c_n\}, \{p_n>0\}, \{q_n\}, \{r_n\} \{s_n\}\) be sequences, and let \( \lambda_1>0 \) be a constant. According to our definition of EOPs (see Def \ref{def:orthonormalEOPs}), we further assume  $\pi_1(z)=0$, and $\pi_{k}=0$ for $k<0$. 
% \rmkH{The sequences \( \{a_n>0\}, \{b_n\}, \{c_n\}, \{p_n>0\}, \{q_n\}, \{r_n\} \{s_n\}\) are? (Note to self:) Discuss this theorem overall. }
Consider the five and seven-term recurrence relations for the polynomials $\pi_k(z)$ respectively
% \footnote{by a slight abuse of notation we use $Q_n$, $R_n$ to denote orthonormal polynomials. However, similar relations hold for orthogonal polynomials with suitable modifed coefficients.}
%\rmkH{Write this statement more clearly.}

\begin{equation}\label{5t}
   - a_{n-1} \pi_n(z) = b_{n-1} \pi_{n-1}(z) + (c_{n-2} - \wp(z)) \pi_{n-2}(z) + b_{n-2} \pi_{n-}(z) + a_{n-3} \pi_{n-4}(z),
\end{equation}
and
\begin{equation}\label{7t}
\begin{aligned}
-p_n \pi_n(z) = q_{n-1} \pi_{n-1}(z) + r_{n-2} \pi_{n-2}(z) + ( s_{n-3}-&\wp'(z)) \pi_{n-3}(z) + r_{n-3} \pi_{n-4}(z)\\& + q_{n-3} \pi_{n-5}(z) + p_{n-3} \pi_{n-6}(z).
\end{aligned}\end{equation}
 If the system is consistent, then there exists a unique positive definite moment functional \( \mathcal{L}: \mathcal{P} \to \mathbb{R} \) such that:
\[
\mathcal{L}(1) = \lambda_1, \quad \mathcal{L}(\pi_n(z) \pi_m(z)) = \lambda_n \delta_{m,n},
\quad \lambda_n>0.\]
\end{theorem}

\begin{proof}
We define the moment functional \( \mathcal{L} \) on the space of polynomials by specifying its action on the orthogonal polynomial sequence \( \{\pi_n(z)\}_{n \geq 0} \) as follows:
\[
\mathcal{L}(1) = \lambda_1, \qquad \mathcal{L}(\pi_n(z)) = 0 \quad \text{for all } n \geq 1.
\]
This definition extends uniquely to all polynomials via linearity.

Let \( n = 2k \) be even, and fix an integer \( i \leq k - 1 \). From the five-term recurrence relation satisfied by the orthogonal polynomials (denoted as equation~\eqref{5t}), we can write:
\[
\wp(x)^i \pi_{2k}(z) = \sum_{j = 2k - 2i}^{2k + 2i} d_j \pi_j(z), \quad \text{for some } d_j \in \mathbb{R}.
\]
Applying \( \mathcal{L} \) to both sides and using the defining property that \( \mathcal{L}(\pi_j(z)) = 0 \) for \( j \geq 1 \), we find:
\[
\mathcal{L}(\wp(x)^i \pi_{2k}(z)) = \sum_{j = 2k - 2i}^{2k + 2i} d_j \mathcal{L}(\pi_j(z)) = 0.
\]

Now consider the expression \( \mathcal{L}(\wp'(z) \wp^i(z) \pi_{2k}(z)) \) with $i\leq k-2$. Using the same recurrence, we write:
\begin{equation}\label{eq:sum_wpprime}
\mathcal{L}(\wp'(z) \wp^i(z) \pi_{2k}(z)) = \sum_{j = 2k - 2i}^{2k + 2i} \mathcal{L}(\wp'(z) d_j \pi_j(z)).
\end{equation}
Each term of the form \( \wp'(z) \pi_j(z) \) can be expressed using the seven-term recurrence relation ~\eqref{7t}. For example, consider the highest-degree term for \( i = k - 2 \), namely:
\[
\wp'(z) \pi_4(z) = \sum_{j = 2}^{7} e_j \pi_j(z), \quad \text{for some } e_j \in \mathbb{R}.
\]
Then,
\[
\mathcal{L}(d_4 \wp'(z) \pi_4(z)) = \sum_{j = 2}^{7} e_j \mathcal{L}(\pi_j(z)) = 0.
\]
As all such terms in~\eqref{eq:sum_wpprime} vanish by the same reasoning, we conclude:
\[
\mathcal{L}(\wp'(z) \wp^i(z) \pi_{2k}(z)) = 0 \quad \text{for all } i \leq k - 2.
\]

A completely analogous argument applies in the case where \( n = 2k + 1 \) is odd. Thus, for any \( m < n \), the inner product induced by \( \mathcal{L} \) satisfies:
\[
\mathcal{L}(\pi_m(z) \pi_n(z)) = 0,
\]
which shows that \( \{\pi_n(z)\} \) is an orthogonal system under \( \mathcal{L} \).

To study the moments, observe that applying \( \mathcal{L} \) to the recurrence for \( \pi_{2n}(z) \), we get:
\[
\mathcal{L}(\wp(z)^n \pi_{2n}(z)) = a_{2n-1} \mathcal{L}(\wp(z)^{n-1} \pi_{2n-2}(z)).
\]
This recurrence allows us to compute moments recursively from lower degrees.

Similarly,
\[
\mathcal{L}(\wp'(z) \wp^n(z) \pi_{2n+3}(z)) = a_{2n+2} \mathcal{L}(\wp'(z) \wp^{n-1}(z) \pi_{2n+1}(z)),
\]


Since the action of \( \mathcal{L} \) vanishes on all polynomials orthogonal to the constant function and satisfies a recursive moment structure on the diagonal, the orthogonality is complete and well-posed. Furthermore, assuming the recurrence coefficients \( a_n, p_n > 0 \), the functional \( \mathcal{L} \) is also positive definite.
\end{proof}

\end{document}
